\clearpage
\fancyhead[L]{\small PALLAS.}
\hypertarget{native-267-lower}{}
\begin{center}
PALLAS.

\rule{0.5\linewidth}{0.4pt}

Monatliche Correspondenz. Band XI. S. 376–383. April 1805.

\rule{0.5\linewidth}{0.4pt}
\end{center}

\emph{Pallas} ließ sich zwar schon seit mehreren Monaten nicht mehr beobachten. Dank des Fleißes von Dr. \textsc{Gauss} können wir unseren Lesern hier aber noch einiges zur Theorie dieses neuen Planeten mitteilen. Das wird vor allem helfen, ihn in der nächsten, wahrscheinlich noch etwas entfernten Sichtbarkeitsperiode wiederzufinden. Nur ein Mangel an Beobachtungsmaterial kann dem unermüdlichen Fleiß von Dr. \textsc{Gauss} und seinem anhaltenden Bemühen um immer genauere berechnete Planetenbahnen Grenzen setzen. Mehrmals äußerte er seine Unzufriedenheit darüber, dass so wenige Beobachtungen der Pallas vorlagen, denn dadurch konnte er die VII. Elemente dieses Planeten nicht weiter berichtigen. Freilich brachte die Beobachtung dieses lichtschwachen Himmelskörpers mancherlei Schwierigkeiten mit sich. Die Beobachtungen von \textsc{Olbers}, \textsc{Oriani}, \textsc{David} und an der hiesigen Sternwarte sind die einzigen, die im ganzen vergangenen Jahr gemacht wurden. Sehr willkommen war ihm daher die am 30. August 1804 in \emph{Brera} beobachtete Opposition der Pallas. Zu diesem Zeitpunkt betrug der mittlere Fehler der VII. Elemente in der Länge $-7'\,28''3$ und in der Breite $+2'\,14''7$. Auf diese Beobachtung und alle verfügbaren früheren Beobachtungen stützte er die folgenden VIII. Elemente der Pallas.
\SourceLeaf{268}
\begin{center}
\begingroup\fontsize{9}{11}\selectfont\setlength{\tabcolsep}{3pt}\renewcommand{\arraystretch}{1.1}
\begin{tabular}{@{}l|l@{}}
\begin{tabular}[t]{@{}lr@{}}
\shortstack[l]{Epoche, Seeberger\\Meridian}  & $\left\{\begin{array}{lr}1803\ .&221^\circ\ 31'\ 23''2\\1804\ .&299\ 58\ 38.1\end{array}\right.$ \\
\shortstack[l]{Tägliche tropische\\Bewegung . . .} & $771''6802$ \\
\shortstack[l]{Aphel 1803\\ . . . \raisebox{.3ex}{.} . . .} & $301^\circ\ 1'\ 44''1$ \\
\end{tabular} &
\begin{tabular}[t]{@{}lr@{}}
\shortstack[l]{Aufsteigender\\Knoten 1803 . . . .} & $172^\circ\ 29'\ 6''8$ \\
\shortstack[l]{Exzentrizität\\ . . . . . . . . .} & $0.246101$ \\
\shortstack[l]{Logarithmus der\\Halbachse . . .} & $0.4417647$ \\
\shortstack[l]{Neigung der\\Bahn . . . . . . .} & $34^\circ\ 37'\ 43''2$ \\
\end{tabular} \\
\end{tabular}\endgroup
\end{center}

\indent Aus diesen Elementen berechnete Dr. \textsc{Gauss} außerdem die folgende Ephemeride für den geozentrischen Lauf der \emph{Pallas} [vom 28. Juli 1805 bis zum 30. April 1806 und gibt die Positionen in derselben Weise an wie die der Ceres für dieselben Jahre, \hyperlink{native-267}{S. 262} dieses Bandes.]

\Needspace{19\baselineskip}
—\quad—\quad—

Dr. \textsc{Gauss} schrieb uns, dass diese Ephemeride noch einmal auf rein elliptischen Elementen beruhe. Er halte es bei Pallas noch nicht für zweckmäßig, die Störungen zu berücksichtigen. Zum einen lassen sich alle bisherigen Beobachtungen noch recht gut durch eine reine Ellipse darstellen. Zum anderen ließ sich die Bahn der Pallas in der kurzen Zeit ihres Erscheinens noch nicht so genau bestimmen, dass man sich von der Entwicklung sämtlicher Störungsgleichungen eine Genauigkeit versprechen dürfte, die diese mühsame Arbeit lohnen könnte. —\quad—\quad—\quad—\quad—\quad— Wäre die Masse der Ceres etwas größer, so könnte sie erhebliche Störungen in der Bahn der Pallas hervorrufen, teils wegen des geringen Abstands zwischen beiden Planeten, teils wegen ihrer fast gleichen mittleren Bewegung.

—\quad—\quad—

Inspektor \textsc{Harding}, dem Dr. \textsc{Gauss} die oben angegebene neue Ephemeride für Pallas sofort mitteilte, arbeitete an der sehr verdienstvollen Aufgabe, eine Sternkarte für ihren Lauf zu entwerfen. Leider werden ihm hier oft Sternpositionen fehlen, da ihm die Histoire céleste in diesen südlichen Zonen weniger hilft, als es diese wertvolle Sammlung sonst tut.

\hfill[\textsc{von Zach.}]
